\section*{Capítulo \ref{ch:g12:work-and-energy} --- Trabalho e energia mecânica}

\begin{solution}{exo:g12:work-and-energy:1}
$E_p = \tfrac12 \times 200 \times 0.050^2 = \qty{0.25}{J}$. Com o
alongamento dobrado: $\times 4$, ou seja, \qty{1.0}{J}. Os segundos
\qty{5.0}{cm} sozinhos custam $1.0 - 0.25 = \qty{0.75}{J}$ --- a \omterm{def:g10:inertia:force}{força}
ali é maior.
\end{solution}

\begin{solution}{exo:g12:work-and-energy:2}
$W = mg(z_A - z_B) = 0.90 \times 9.81 \times (12 - 30) \approx
\qty{-1.6e2}{J}$, seja qual for a rota --- o mesmo pela subida reta. Na
ida e volta: \qty{0}{J} (\omterm{def:g12:work-and-energy:conservative}{força conservativa}).
\end{solution}

\begin{solution}{exo:g12:work-and-energy:3}
(a), (b) \omterm{def:g12:work-and-energy:conservative}{conservativas}: \omterm{def:g11:work-of-force:work}{trabalho} fixado pelos extremos (desnível de
altitude; alongamento). (c), (e) \omterm{def:g12:work-and-energy:conservative}{não conservativas}: $-fL$ e o arrasto
cobram por \omterm{def:g10:orders-of-magnitude:unit}{metro} percorrido. (d) sem \omterm{def:g11:work-of-force:work}{trabalho}: sempre perpendicular ao
movimento.
\end{solution}

\begin{solution}{exo:g12:work-and-energy:4}
$h = L(1 - \cos\qty{60}{\degree}) = \qty{1.0}{m}$;
$v = \sqrt{2 \times 9.81 \times 1.0} \approx \qty{4.4}{m/s}$. A tensão é
perpendicular ao movimento a cada instante: \omterm{def:g11:work-of-force:power}{potência} nula, \omterm{def:g11:work-of-force:work}{trabalho} nulo.
\end{solution}

\begin{solution}{exo:g12:work-and-energy:5}
$F = P/v = 250/9.0 \approx \qty{28}{N}$. A \qty{5.0}{m/s}:
$F = 250/5.0 = \qty{50}{N}$ (na maior parte, a componente do \omterm{def:g10:universal-gravitation:weight}{peso} ao
longo da subida).
\end{solution}

\begin{solution}{exo:g12:work-and-energy:6}
$E_p = \tfrac12 \times 450 \times 0.080^2 = \qty{1.44}{J}$;
$v = \sqrt{2 \times 1.44/0.020} = \qty{12}{m/s}$;
$h = E_p/mg = 1.44/(0.020 \times 9.81) \approx \qty{7.3}{m}$.
\end{solution}

\begin{solution}{exo:g12:work-and-energy:7}
Trapézios (passos de \qty{2.0}{cm}): $(0.03 + 0.10 + 0.19 + 0.30) =
\qty{0.62}{J}$. Modelo de reta: $k = 18.0/0.080 = \qty{225}{N/m}$,
$\tfrac12 kx^2 = \qty{0.72}{J}$. A curva medida afunda abaixo da reta
para $x$ pequeno, de modo que a área verdadeira é menor.
\end{solution}

\begin{solution}{exo:g12:work-and-energy:8}
$\Delta E_m = \tfrac12 \times 55 \times 9.0^2 - 55 \times 9.81 \times
8.0 = 2228 - 4316 \approx \qty{-2.1e3}{J}$;
$f = 2089/25 \approx \qty{84}{N}$. O \omterm{def:g10:universal-gravitation:weight}{peso} é conservativo e a \omterm{def:g12:newtons-laws:friction}{força normal}
não trabalha: entram apenas o desnível e o \emph{comprimento} do
caminho, e não a forma dele.
\end{solution}

\begin{solution}{exo:g12:work-and-energy:9}
$H_{\min} = \tfrac52 R = \qty{0.50}{m}$. De \qty{60}{cm}:
$v_{\text{alto}} = \sqrt{2 \times 9.81 \times (0.60 - 0.40)} \approx
\qty{2.0}{m/s}$, e $v_{\text{alto}}^2/(gR) = 2.0$ --- o dobro do mínimo.
O \omterm{def:g10:inertia:inventory}{atrito} raspa \omterm{def:g10:energy-conservation:energy}{energia} o tempo todo, e por isso as pistas reais precisam
de altura extra.
\end{solution}

\begin{solution}{exo:g12:work-and-energy:10}
\omterm{def:g12:work-and-energy:diagram}{Pontos de retorno}: $8.0\,x^2 = 2.0$, $x = \pm\qty{0.50}{m}$. Velocidade
máxima em $x = 0$: $v = \sqrt{2 \times 2.0/0.100} \approx \qty{6.3}{m/s}$.
Uma oscilação de vaivém num poço parabólico --- a $E_p$ do \omterm{prop:g12:oscillators-and-time:spring-period}{oscilador
massa--mola} tem exatamente essa forma.
\end{solution}

\begin{solution}{exo:g12:work-and-energy:11}
$\tfrac12 kx^2 = \tfrac12 mv^2$:
$x = v\sqrt{m/k} = 2.0\sqrt{1200/\num{6.0e5}} \approx \qty{8.9}{cm}$. A
\qty{4.0}{m/s}: \qty{18}{cm} --- o dobro, e não o quádruplo: a \omterm{def:g10:energy-conservation:energy}{energia}
quadruplica, mas $E_p$ cresce como $x^2$, de modo que $x$ apenas dobra.
\end{solution}

\begin{solution}{exo:g12:work-and-energy:12}
(a) $E_m = \qty{5.0}{J} < \qty{6.0}{J}$: ele não consegue passar; para na
encosta onde $E_p = \qty{5.0}{J}$ e volta rolando. (b) No alto:
$E_k = \qty{1.0}{J}$, $v = \sqrt{2 \times 1.0/0.50} = \qty{2.0}{m/s}$; no
vale: $E_k = \qty{5.0}{J}$, $v \approx \qty{4.5}{m/s}$. (c) $x = 0$:
instável (máximo); $x = \qty{1.5}{m}$: estável (mínimo).
\end{solution}

\begin{solution}{exo:g12:work-and-energy:13}
$\tfrac12 mv^2 = mgR_E$:
$v = \sqrt{2 \times 9.81 \times \num{6.37e6}} \approx \qty{1.12e4}{m/s}
= \qty{11.2}{km/s}$. Todo termo é proporcional a $m$: sonda e seixo por
igual. Logo abaixo dela, o lançamento ainda sobe enormemente longe,
encontra um \omterm{def:g12:work-and-energy:diagram}{ponto de retorno} e cai de volta: preso.
\end{solution}

\begin{solution}{exo:g12:work-and-energy:14}
$E_m = \tfrac12 kA^2 = \tfrac12 \times 40 \times 0.060^2 =
\qty{7.2e-2}{J}$; $v_{\max} = \sqrt{2E_m/m} =
\sqrt{2 \times 0.072/0.250} \approx \qty{0.76}{m/s}$; $E_k = E_p$ onde
$\tfrac12 kx^2 = E_m/2$: $x = \pm A/\sqrt 2 \approx \pm\qty{4.2}{cm}$.
\end{solution}

\begin{solution}{exo:g12:work-and-energy:15}
$h = v^2/2g = 9.5^2/19.62 \approx \qty{4.6}{m}$; sarrafo $\approx
4.6 + 1.0 = \qty{5.6}{m}$. O recorde é mais alto: os saltadores
acrescentam \omterm{def:g11:work-of-force:work}{trabalho} com braços e pernas sobre a vara que verga. Mas a
estimativa energética fixa a escala --- nenhum salto de velocista
chegará a \qty{10}{m}.
\end{solution}

\begin{solution}{pb:g12:work-and-energy:1}
\textbf{1.} $v = \sqrt{2gH} = \sqrt{2 \times 9.81 \times 45} \approx
\qty{29.7}{m/s}$: \omterm{def:g10:universal-gravitation:weight}{peso} conservativo, \omterm{def:g12:newtons-laws:friction}{força normal} sem \omterm{def:g11:work-of-force:work}{trabalho} --- o
perfil sai da conta.

\textbf{2.} Portão: $E_m = mgH = 70 \times 9.81 \times 45 \approx
\qty{3.09e4}{J}$. Saída: $E_m = \tfrac12 \times 70 \times 26^2 \approx
\qty{2.37e4}{J}$.

\textbf{3.} $\Delta E_m \approx \qty{-7.2e3}{J}$; o \omterm{thm:g12:work-and-energy:emtheorem}{teorema da energia
mecânica} cobra isso das \omterm{def:g12:work-and-energy:conservative}{forças não conservativas}: o \omterm{def:g10:inertia:inventory}{atrito} da neve e o
arrasto do ar.

\textbf{4.} $f = \num{7.2e3}/90 \approx \qty{80}{N}$.

\textbf{5.} $\num{7.2e3}/\num{3.09e4} \approx 23\%$. A cera abaixa $f$ e
o agachamento abaixa o arrasto: os dois encolhem $fL$, o único termo
negociável.

\textbf{6.} \omterm{def:g12:projectile-motion:freefall}{Queda livre} com velocidade inicial horizontal --- um
lançamento de \omterm{def:g12:projectile-motion:projectile}{projétil}. $v_x = \qty{26}{m/s}$;
$v_y = \sqrt{2 \times 9.81 \times 40} \approx \qty{28.0}{m/s}$.

\textbf{7.} $v = \sqrt{26^2 + 28.0^2} = \sqrt{1461} \approx
\qty{38.2}{m/s}$; pela \omterm{def:g10:energy-conservation:energy}{energia}: $v = \sqrt{v_0^2 + 2gh} =
\sqrt{676 + 785}$ --- idêntico. Cerca de \qty{138}{km/h}.

\textbf{8.} $t = v_y/g = 28.0/9.81 \approx \qty{2.9}{s}$;
$d = v_0 t \approx \qty{74}{m}$.

\textbf{9.} Não. A sustentação é perpendicular à velocidade (não
trabalha) e o arrasto se opõe a ela (\omterm{def:g11:work-of-force:work}{trabalho} negativo):
$\Delta E_m \leq 0$, de modo que a velocidade real de aterrissagem só
pode ser \emph{menor} --- o ar estica o voo, ele não o alimenta.

\textbf{10.} $\alpha = \arctan(28.0/26) \approx \qty{47}{\degree}$ abaixo
da horizontal.

\textbf{11.} Ângulo com a encosta: $47 - 35 = \qty{12}{\degree}$.
$v_\parallel = 38.2\cos\qty{12}{\degree} \approx \qty{37.4}{m/s}$;
$v_\perp = 38.2\sin\qty{12}{\degree} \approx \qty{7.9}{m/s}$.

\textbf{12.} Absorvida: $\tfrac12 \times 70 \times 7.9^2 \approx
\qty{2.2e3}{J}$, contra $\tfrac12 \times 70 \times 38.2^2 \approx
\qty{5.1e4}{J}$ no total: cerca de $4\%$.

\textbf{13.} $h_{\text{eq}} = 7.9^2/19.62 \approx \qty{3.2}{m}$ --- um
salto ousado de um muro. No plano: $38.2^2/19.62 \approx \qty{74}{m}$ ---
insobrevivível.

\textbf{14.} Uma encosta paralela à \omterm{def:g10:relative-motion:trajectory}{trajetória} do voo mantém $v_\perp$
pequena, de modo que as pernas absorvem \omterm{def:g10:orders-of-magnitude:unit}{metros}, e não a montanha inteira
de \omterm{def:g11:mechanical-energy:kinetic}{energia cinética}.

\textbf{15.} $H'_0 = v_0'^2/2g = 576/19.62 \approx \qty{29.4}{m}$.

\textbf{16.} $mgH' - f\,\dfrac{H'}{\sin\qty{35}{\degree}} =
\tfrac12 m v_0'^2$.

\textbf{17.} $H' = \dfrac{\tfrac12 \times 70 \times 576}
{686.7 - 80/0.574} = \dfrac{\num{20160}}{547} \approx \qty{37}{m}$.

\textbf{18.} $36.8/29.4 \approx 1.25$: uma sobretaxa de $25\%$.

\textbf{19.} Para \qty{60}{kg}: $H' = \num{17280}/(588.6 - 139.5)
\approx \qty{38.5}{m}$. A saltadora \emph{mais leve} precisa da torre
mais alta: a conta do \omterm{def:g10:inertia:inventory}{atrito} $fL$ é a mesma, mas o orçamento de \omterm{def:g10:energy-conservation:energy}{energia}
$mgH'$ dela é menor.

\textbf{20.} Anuncie uma torre de \qty{37}{m} para \qty{24}{m/s} de
saída --- e confesse que um quarto dela, uns \qty{7}{m} de concreto, é a
sobretaxa que o \omterm{def:g10:inertia:inventory}{atrito} cobra do sonho sem \omterm{def:g10:inertia:inventory}{atrito}.
\end{solution}
